% !TEX root = ../MasterThesis.tex
\normallinespacing

\chapter{Benchmarks}
\label{chap:benchmarks}

In this project we compare the strengths and weaknesses of a range of
certificate verification methods. The suite therefore spans a broad range of
system types, drawn from the literature to expose where each method succeeds and
fails. By identifying structural features of systems that make verification easy
or hard, we can separate each method along interpretable axes.

We select systems from the formal-methods and control literature so that the
suite spans three orthogonal axes of difficulty:

\begin{itemize}
  \item \textbf{Dynamics class.} From linear, through polynomial, to
        transcendental (e.g.\ trigonometric) dynamics. This axis directly controls
        which symbolic engines can encode the one-step map \emph{exactly}.
  \item \textbf{Noise model.} We assume that the noise distribution is bounded for each system.
        We include systems ranging from discrete to continuous additive noise variables,
        and from internal, state-dependent noise to external additive noise.
  \item \textbf{Dimension.} The linear families scale to arbitrary dimension,
        allowing us to explore how the curse of dimensionality affects each method.
\end{itemize}

A second selection principle is the \textbf{certificate class} each system
admits. Some systems have a known, analytic ground-truth certificate (a
quadratic $x^\top P x$, an $L_1$ piecewise-linear function, or a physical
potential), which lets us measure a verifier against a \emph{known-good} answer;
others have no closed-form certificate and the certificate must be learned,
exercising the full CEGIS loop.

Each benchmark instantiates the discrete-time stochastic system
\eqref{eq:stochastic-system}: its dynamics and noise model fix the transition
kernel $P(\cdot \mid x)$ of Section~\ref{subsec:kernels}, and with it the
one-step expectation inside the drift \eqref{eq:drift-def} that every verifier
of Chapter~\ref{chap:verification} must bound. The selection axes are therefore
precisely the levers that control how hard that expectation is to work with:
the dynamics class governs whether the one-step map can be encoded symbolically, the
noise model governs whether the expectation is an exact finite sum or must be
soundly over-approximated, and the dimension scales both.

\begin{table}[h]
\centering
\caption{The benchmark suite across the three selection axes.}
\label{tab:bench-suite}
\begin{tabular}{llll}
\hline
\textbf{System} & \textbf{Dynamics} & \textbf{Noise} & \textbf{Analytic Certificate} \\
\hline
$\mathrm{LinSwitch}$    & PWL / jump     & multiplicative (switch) & $L_1$ PWL (exact)  \\
$\mathrm{LinStoch}$   & linear         & additive bounded        & quadratic $x^\top P x$ \\
$\mathrm{DoubleWell}$        & polynomial     & additive bounded        & potential\\
$\mathrm{Pendulum}$     & transcendental & additive triangular     & none \\
\hline
\end{tabular}
\end{table}

Phase portraits of all four 2D systems (the expected drift field with sample
trajectories overlaid) are collected in Appendix~\ref{app:benchmark-portraits}.

% =====================================================================
\section{Linear Systems}
\label{sec:bench-linear}

% ---------------------------------------------------------------------
\subsection{$\mathrm{LinSwitch}$}
\label{sec:bench-switched}

The randomly switched linear system is a canonical model in the stochastic
stability literature, where at each step the dynamics matrix is selected by an
i.i.d.\ Bernoulli switch:
\[
  x_{t+1} = A_{\sigma_t} x_t, \qquad \sigma_t \sim \text{Bernoulli}(p).
\]
We instantiate the randomly switched systems framework of Chatterjee \&
Pal~\cite{chatterjee2009excursion}, who give conditions for almost-sure and
$L_1$ stability of such systems. Our 2D instance uses a pair of matrices chosen
to satisfy their contractivity conditions, being jointly contractive in the
$L_1$ norm. Randomly switched systems model phenomena such as mode-switching
control, packet loss over lossy channels, and Markov-jump linear dynamics. A
phase portrait of the 2D instance is given in Figure~\ref{fig:phase-linswitch}.

This is the \emph{simplest} stochastic obstruction in the suite and our
soundness baseline. The switch is an \emph{internal, discrete} source of
randomness, so it is encoded as a small set of exact finite branches rather than
an over-approximated continuous noise box, leaving no relaxation anywhere in the
pipeline. Every engine handles it, which makes it the control case against which
we measure. The $L_1$-contractivity gives an \emph{exact piecewise-linear}
ground-truth certificate $\cert(x) = |x_0| + |x_1|$, so the verification
conditions hold with \emph{zero relaxation gap} in both SMT and MILP, a property
no other system in the suite shares. The construction also generalises. A
random-construction variant draws two matrices at a target spectral norm $<1$ in
any dimension, for which $\cert(x) = \lVert x \rVert_2$ is valid, letting us
probe dimensional scaling.


% ---------------------------------------------------------------------
\subsection{$\mathrm{LinStoch}$}
\label{sec:bench-linstoch}

The additive-noise linear system,
\[
  x_{t+1} = A x_t + w_t, \qquad w_t \sim \text{bounded zero-mean noise},
\]
with $A$ Schur-stable, is \emph{the} canonical 2D example of the
stochastic-supermartingale literature~\cite{lechner2022stability, badings2025logRASM}. The deterministic part is
contractive, so the only obstruction to stability is the noise floor.
Figure~\ref{fig:phase-linstoch} shows the resulting inward spiral to the
noise-floor neighbourhood.

It is the additive-noise counterpart to the $\mathrm{LinSwitch}$ and isolates the
\emph{sound treatment of continuous additive noise}. The noise distribution is
over-approximated by a discretisation into cells so that the symbolic engines
soundly bound the expected drift. It is also the natural scaling testbed, since
$A$ can be constructed at any dimension and target spectral norm.

It admits an \emph{analytic} quadratic ground-truth certificate $\cert(x) = x^\top P
x$ with $P$ the solution of the discrete Lyapunov equation $A^\top P A - P = -Q$
($Q = I$). The expected drift then has the closed form
\[
  \expe[\cert(x')] - \cert(x) = -x^\top Q x + \operatorname{tr}(P\Sigma_w),
\]
which is $\le 0$ exactly outside a ball of radius
$r^\star = \sqrt{\operatorname{tr}(P\Sigma_w)/\lambda_{\min}(Q)}$. The
equilibrium radius is therefore \emph{derived from the noise level}, so $\cert$ is a
genuine supermartingale on $\text{domain}\setminus\text{equilibrium}$ (the
textbook ``a.s.\ stable to a neighbourhood'').


% ---------------------------------------------------------------------
\subsection{$\mathrm{DoubleWell}$}
\label{sec:bench-doublewell}

A 2D stochastic gradient-flow system on a tilted double-well potential,
\[
  x_{k+1} = x_k - \eta\,\nabla U(x_k) + w_k, \qquad
  U(x,y) = \tfrac14 (x^2 - 1)^2 + \tfrac12 y^2 - 0.3\,x .
\]
The double-well is the canonical model of metastability and noise-driven escape
over a potential barrier, originating with Kramers~\cite{kramers1940brownian} and
recurring throughout the Lyapunov-certificate literature as a bistable benchmark.
Figure~\ref{fig:phase-doublewell} shows the flow around the off-origin attractor.

It is the first \emph{nonlinear} (polynomial) system, while still
admitting an exact symbolic encoding (the cubic gradient is encoded via an
auxiliary squared variable in Z3/Gurobi). It tests whether a verifier handles a
non-convex value landscape and a non-origin equilibrium. The potential is
\emph{bistable}. The tilt $-0.3x$ breaks the symmetry so one well is the true
attractor, but the second (metastable) well is a basin the certificate must rule
out. The potential $U$ itself is the ground-truth Lyapunov function, giving an
interpretable, physically meaningful $\cert$. The equilibrium is off-origin (at
$\approx(1,0)$), unlike every other system.


\subsection{$\mathrm{Pendulum}$}
\label{sec:bench-pendulum-lqr}

The inverted pendulum is the archetypal control benchmark. We take the FOSSIL
\texttt{SineModel} pendulum with its verified LQR feedback gain $K = [[7.21,
1.34],[1.34, 0.33]]$~\cite{edwards2024fossil2, abate2021fossil},
folded into autonomous closed-loop dynamics, with additive triangular process
noise:
\[
  \theta' = \theta + h(\dot\theta + u_1), \qquad
  \dot\theta' = \dot\theta + h\!\left(u_2 + \tfrac{mgL\sin\theta - b\dot\theta}{mL^2}\right),
  \quad u = -Kx .
\]

This is the first \emph{transcendental} system. The $\sin\theta$ term means the
SMT and MILP engines cannot encode the map and must abstain, so only autoLiRPA
(which relaxes $\sin$) and the sampling/MAB engines apply. It also represents
dynamics found in the control literature, in which a controller returns actions
to influence a trajectory. Here $u$ is the controller and represents a torque
applied to the pendulum, and the equilibrium is its upright position, which would
otherwise be unstable if the controller was poor. The closed loop is Schur-stable but \emph{non-normal}
($\lVert A_d \rVert_2 > 1$, transient growth), so neither a box nor a ball is
forward invariant, and the domain is the Lyapunov sublevel ellipsoid of the
linearisation. Nonlinear dynamics and a controller add complexity beyond that of
the previous systems. Figure~\ref{fig:phase-pendulum} shows the closed-loop
spiral over the ellipsoidal domain.


\newpage
